\section{System Overview}
The system consists of three modules:
\begin{itemize}[leftmargin=*]
  \item Android endpoint app (\texttt{VpnService}, flow capture, on-device scoring, UI).
  \item Backend adapter (authenticated ingest, queue/retry, triage, policy APIs, optional SIEM forwarding).
  \item ML pipeline (feature engineering, training, evaluation, threshold calibration, model export).
\end{itemize}
The architectural goal is to maximize local autonomy while still supporting optional centralized
analysis and reproducible experimentation. Detection must remain possible without backend
connectivity. The same feature and policy definitions must also support export, privacy-view
derivation, leakage benchmarking, and model retraining.

\begin{figure}[htbp]
  \centering
  \fbox{%
    \begin{minipage}{0.92\textwidth}
      \centering
      \textbf{Android endpoint}\\
      \small \texttt{VpnService} capture \(\rightarrow\) flow metadata \(\rightarrow\) local scoring
      \(\rightarrow\) privacy-tiered export

      \vspace{0.6em}
      \(\Downarrow\)

      \vspace{0.6em}
      \textbf{Backend adapter}\\
      \small authenticated ingest \(\rightarrow\) queue/retry \(\rightarrow\) triage and policy
      \(\rightarrow\) optional SIEM/Wazuh forwarding

      \vspace{0.6em}
      \(\Downarrow\)

      \vspace{0.6em}
      \textbf{ML and evidence pipeline}\\
      \small feature windows \(\rightarrow\) model training \(\rightarrow\) privacy leakage audits
      \(\rightarrow\) archived thesis reports
    \end{minipage}%
  }
  \caption{High-level MANTA data path from endpoint capture to evaluation evidence.}
  \label{fig:manta-architecture}
\end{figure}

\section{Data Flow}
The runtime data path is:
\begin{enumerate}[leftmargin=*]
  \item Packet capture via the \texttt{VpnService} TUN interface.
  \item Userspace forwarding to preserve normal network connectivity.
  \item Packet parsing and app attribution to flow records.
  \item Per-app rolling window aggregation. The default conceptual window is 60 s, with adaptive
  30 s and 300 s modes for dense or sparse/service traffic.
  \item Local anomaly scoring through statistical, multivariate, sequence, exported-tree, hybrid,
  TFLite-compatible, remote-assisted, and fusion modes.
  \item Response-score adjustment, severity resolution, stability guards, false-positive budget
  checks, and alert-evidence policy.
  \item Local persistence of flows, windows, and alerts in Room.
  \item Optional export of selected flow and alert events to backend APIs.
\end{enumerate}

Policy synchronization is decoupled through background workers. Threshold updates and
retention/export changes can therefore be applied without requiring a new app build.

\section{Threat Model and Assumptions}
The primary defensive goal is to detect suspicious encrypted traffic behavior. Examples include
beaconing, unusual destination changes, remote-control style communication,
data-exfiltration-like volume shifts, and malware-associated anomalies. In the security framing from
Chapter 1, these behaviors matter because they may indicate compromise of confidentiality,
integrity, availability, or security accountability. The design assumes that payloads remain
encrypted and unavailable for inspection.

The thesis considers three privacy-relevant adversaries:
\begin{itemize}[leftmargin=*]
  \item a \textbf{data consumer} who receives an exported or derived representation and attempts to
  infer application or behavior properties from that release;
  \item a \textbf{passive observer} who watches encrypted traffic on the network path and performs
  fingerprinting from sequence metadata alone;
  \item a \textbf{central coordinator} in collaborative training who may receive learned updates or
  reduced latent features rather than raw full-feature traffic.
\end{itemize}

The model does \emph{not} assume root compromise of the endpoint OS, payload decryption on the wire,
or strong traffic-shaping defenses such as dummy-packet padding.

Trust boundaries are defined as follows:
\begin{itemize}[leftmargin=*]
  \item \textbf{Trusted}: endpoint app process, local encrypted settings, backend adapter process.
  \item \textbf{Partially trusted}: network transport path, which is TLS-protected but still exposed
  to outages and passive observation.
  \item \textbf{Untrusted}: remote application peers and public network infrastructure.
\end{itemize}

\section{Privacy Model}
MANTA separates two privacy and leakage surfaces:
\begin{itemize}[leftmargin=*]
  \item \textbf{Release privacy}: what can be inferred from the feature representation that is
  intentionally exported, shared, stored, or used as the privacy-reduced model input.
  \item \textbf{Observer leakage}: what can still be inferred by a passive eavesdropper from the raw
  encrypted traffic sequence itself.
\end{itemize}

The distinction drives the design. Feature hashing, identifier suppression, bucketization,
privacy-student training, and federated latent sharing act on the released feature representation.
A passive observer sees the original encrypted sequence before those transformations occur.

\begin{figure}[htbp]
  \centering
  \fbox{%
    \begin{minipage}{0.92\textwidth}
      \textbf{Release privacy:} controls and audits the feature representation intentionally
      exported by MANTA. This surface is affected by hashing, removal of direct identifiers,
      bucketization, and privacy-student learning.

      \vspace{0.7em}
      \textbf{Observer leakage:} concerns an on-path observer who sees encrypted-flow timing,
      volume, direction, and sequence patterns before MANTA applies export-time reductions.

      \vspace{0.7em}
      \textbf{Boundary:} feature coarsening can reduce release leakage, but it cannot hide the
      original traffic pattern from an observer on the network path.
    \end{minipage}%
  }
  \caption{Difference between release privacy and observer leakage in MANTA.}
  \label{fig:privacy-surfaces}
\end{figure}

\section{Telemetry-Reduction Design Decisions}
The system applies data minimization and disclosure control at multiple points:
\begin{itemize}[leftmargin=*]
  \item metadata-only flow capture with no payload retention,
  \item hashed or removed direct identifiers in non-\texttt{off} privacy views,
  \item coarsened window features in the \texttt{medium} and \texttt{strict} release views,
  \item privacy-student training to learn a lower-leakage latent representation,
  \item federated-style training over privacy-student latents rather than raw full-feature tables.
\end{itemize}

These mechanisms reduce the representation that MANTA releases, but they are not formal privacy
guarantees. The evaluation measures how much leakage remains after each transformation.

\section{Deployment and Evaluation Boundaries}
The design is anomaly-first rather than anonymity-first. MANTA does not implement traffic-shaping
defenses such as padding, dummy traffic, or timing obfuscation. Its privacy mechanisms operate on
the representation released by the endpoint or backend pipeline. The observer benchmark measures
what remains visible on the network path. In other words, MANTA controls exported telemetry; it does
not implement transport-level anonymity defenses.
